\chapter{CONCLUSION}

This project developed a full-stack food recognition and nutrition-assistance system for Indian meals. The system brings together a YOLO-based detector, a curated 72-class dataset, a SQLite nutrition database, a FastAPI backend, and a Next.js frontend. In the implemented workflow, a user uploads a food image, views detected foods with bounding boxes, checks macro values, adjusts serving assumptions, and receives dietary suggestions based on personal goals and health conditions.

The main objectives were achieved. Five YOLO-format source datasets were merged into a canonical Indian food dataset. The build process normalized raw labels, exported balanced splits, and produced machine-readable artifacts. Visual class verification confirmed that all 72 class names have corresponding image samples and that the class list is aligned. This was important because the dataset had to support model training, nutrition mapping, and frontend display. The nutrition pipeline produced 1,010 food records, including five verified supplemental rows, and created 103 total YOLO-to-database mappings. For the expanded dataset, all 72 classes now have nutrition mappings.

The final YOLO11s detector was trained and evaluated on the expanded dataset. The best exported model achieved 0.830 mAP@50 and 0.692 mAP@50:95 on validation, and 0.820 mAP@50 and 0.680 mAP@50:95 on the held-out test split. The backend and frontend were implemented as a working local prototype, not only as a model experiment. The backend includes image analysis, nutrition validation, macro calculation, goal, health-condition, and recommendation endpoints. The frontend handles image upload, detection visualization, macro display, serving adjustment, and recommendation review.

The work shows that food recognition becomes more useful when it is connected to nutrition databases and user-aware recommendation logic. Dataset normalization and nutrition mapping are not minor details in Indian food systems; they decide whether the output can be trusted. The prototype still has limitations in portion estimation, annotation quality, class confusion among similar dishes, and nutrition source consistency. Even with these limitations, it gives a working foundation for a controlled recommendation model, improved portion handling, mobile use, long-term tracking, and stronger safety review.
